%%
%% CSD Assignment 1 - Final Technical Report
%% AI-based Village Pond Planning System
%% Based on the ACM "acmart" manuscript (single-column) style.
%%
% The libertine font package acmart prefers is not installed in this build
% environment; microtype's font-expansion feature then crashes against the
% Computer-Modern fallback (bitmap fonts at some sizes aren't "scalable").
% Disabling expansion (protrusion/kerning still apply) avoids that crash
% without needing extra font packages.
\PassOptionsToPackage{expansion=false}{microtype}
\documentclass[manuscript,screen]{acmart}

\AtBeginDocument{%
  \providecommand\BibTeX{{Bib\TeX}}}

\settopmatter{printacmref=false}
\renewcommand\footnotetextcopyrightpermission[1]{}
\pagestyle{plain}

\begin{document}

%% =================== TITLE & AUTHOR BLOCK =========================
\title{AI-based Village Pond Planning System}
\subtitle{CSD Assignment 1 -- Final Technical Report}

\author{Paritosh Lahre}
\affiliation{%
  \institution{Indian Institute of Technology Bhilai --- Student ID: 12341550}
  \city{Bhilai}
  \country{India}}
\email{paritoshl@iitbhilai.ac.in}

\renewcommand{\shortauthors}{Lahre}

%% =================== ABSTRACT ======================================
\begin{abstract}
Rural water conservation in India often depends on identifying suitable
sites for small check-dams and percolation ponds, a task traditionally
done by manual site inspection with no systematic terrain, catchment, or
rainfall analysis. This project implements a full-stack, terrain-aware
decision-support system that recommends pond locations from either an
uploaded contour map (KML/KMZ) or a user-drawn area on an interactive map.
The backend reconstructs a digital elevation model, hydrologically
conditions it (Priority-Flood depression filling), computes D8 flow
direction and accumulation, delineates catchments via upstream graph
search, and scores candidate sites on five terrain/hydrology factors. Real
historical rainfall (Open-Meteo) drives a Rational-Method runoff estimate,
which in turn sizes a planning-level pond (depth, surface area, storage
volume). A land-use constraint filter, built on live OpenStreetMap data via
the Overpass API, rejects candidates that fall on existing buildings,
roads, rivers, or water bodies. All of this is exposed through a React and
Leaflet frontend with a dark/light theme and both input modes, engineered
to degrade gracefully rather than fail outright when any of its five
external data dependencies is unavailable -- a design choice validated
repeatedly against real, confirmed failures (DNS resolution errors,
connection resets, and undocumented rate limits from more than one
provider) observed on the actual deployment network. On a representative
$\sim$4\,km\textsuperscript{2} village contour dataset, the system
evaluated 187 candidate seeds, accepted 5 after terrain/catchment/land-use
filtering, and recommended a site with an estimated 21{,}820\,m\textsuperscript{3}
annual collectable runoff and 17{,}456\,m\textsuperscript{3} planned
storage, computed end-to-end in 1.2--7.2 seconds. The map-selected-area
path was verified end-to-end against live external APIs under real,
observed network failures, with a bounded worst-case wall-clock time kept
safely under the client's request timeout. The system is validated by 70
automated tests and extensive live testing against real (not mocked)
external services.
\end{abstract}

\keywords{Village Pond Planning, Geospatial Analysis, Rainwater
  Harvesting, Catchment Estimation, Web GIS, REST APIs}

\maketitle

%% =====================================================================
\section{Introduction}
\label{sec:intro}

Rural India loses a large fraction of monsoon rainfall to unmanaged
runoff. Small ponds and check-dams built at hydrologically correct
locations -- natural depressions or flow-convergence points with adequate
catchment -- can capture this water cheaply, but choosing that location
today relies on manual inspection with no systematic terrain, catchment,
or rainfall analysis. This project builds a web application that
automates that decision: given either a contour map or a map-drawn area,
it reconstructs terrain, analyzes drainage, scores candidate sites,
filters out sites that would conflict with existing land use, estimates
expected water volume from real rainfall data, and presents all of this
overlaid on an interactive map. The rest of this report covers the
requirements (\S\ref{sec:requirements}), architecture
(\S\ref{sec:architecture}), the terrain/hydrology methodology
(\S\ref{sec:methodology}), implementation details
(\S\ref{sec:implementation}), how core Computer System Design themes were
applied (\S\ref{sec:csd-themes}), results (\S\ref{sec:results}), and an
honest discussion of limitations (\S\ref{sec:discussion}).

\subsection{Motivation}
Manual site selection is hard for three concrete reasons this project
directly addresses. First, terrain complexity: a contributing catchment
is not visible by eye, it requires flow-direction and flow-accumulation
analysis across an entire elevation surface. Second, there is no
centralized, per-site rainfall record -- village-level historical
rainfall statistics are not something a field surveyor can look up on the
spot. Third, there is no systematic way to cross-check a promising
terrain location against existing buildings, roads, or rivers without a
separate land-survey pass. Automating these three lets a non-specialist
get a defensible, explainable recommendation in seconds instead of a
multi-day manual survey.

\subsection{Scope of the Project}
\textbf{In scope}: two input modes (upload a KML/KMZ contour map, or draw
an arbitrary area directly on a map, no file needed); terrain
reconstruction, hydrological conditioning, flow analysis, and
multi-candidate catchment delineation from either input; multi-factor
candidate scoring and ranking with human-readable reasoning for the top
recommendation; real historical rainfall retrieval and Rational-Method
runoff estimation; planning-level pond sizing (depth, surface area,
storage volume); a land-use constraint filter (buildings, roads, rivers,
water bodies, power lines) using live OpenStreetMap data; an interactive
map frontend showing every layer above, with graceful degradation
displayed to the user (not hidden) when any external data source is
temporarily unavailable.

\textbf{Out of scope}: structural/civil engineering design of the pond
embankment or spillway -- outputs are planning-level estimates, not
construction drawings; soil type, infiltration, or land-use-based
refinement of the runoff coefficient -- a single configurable constant is
used; verifying land ownership; multi-user authentication or per-user
data isolation; persistent storage beyond a single server process's
lifetime (\S\ref{sec:implementation}); and guaranteeing live external API
uptime -- the system is explicitly designed to degrade and report
``unavailable'' rather than assume any of its five external dependencies
is always reachable, because in practice every one of them was observed
to fail at least once during development.

%% =====================================================================
\section{Problem Statement and Requirements}
\label{sec:requirements}

\begin{table}[h]
  \caption{Functional requirements and where they are implemented}
  \label{tab:requirements}
  \begin{tabular}{p{3.6cm}p{3.6cm}}
    \toprule
    Requirement & Implemented in \\
    \midrule
    Satellite imagery display & \texttt{MapView.jsx} (Esri World Imagery layer) \\
    Contour map visualization & \texttt{kml\_service.py} $\to$ GeoJSON $\to$ \texttt{MapView.jsx} \\
    Available-land identification & \texttt{candidate\_service.py} + \texttt{landuse\_service.py} \\
    Catchment area estimation & \texttt{catchment\_service.py}; flow direction/accumulation and watershed algorithms in \texttt{algorithms/} \\
    Historical rainfall query & \texttt{rainfall\_service.py} (Open-Meteo) \\
    Runoff volume estimation & \texttt{runoff\_service.py} (Rational Method) \\
    Pond depth / storage recommendation & \texttt{pond\_service.py} \\
    Combined overlay / results view & \texttt{MapView.jsx} + \texttt{ResultsPanel.jsx} \\
    \bottomrule
  \end{tabular}
\end{table}

\subsection{Non-Functional Requirements}

\textbf{Performance}: contour-upload analysis completes in under 10
seconds end-to-end for a village-scale ($\sim$4\,km\textsuperscript{2})
dataset -- observed 1.17--7.22\,s pipeline time across multiple real runs
against live external APIs. Map-selected-area analysis is bounded to a
worst case of roughly 75--90 seconds even under significant external-API
degradation, via a hard 45-second elevation-fetch deadline and a
25-second land-use-fetch deadline that runs concurrently with it
(\S\ref{sec:methodology}, \S\ref{sec:discussion}) -- this bound was
discovered necessary after a real production log showed an unbounded,
multi-minute worst case before these deadlines existed.

\textbf{Concurrency/Scalability}: the backend's event loop remains
responsive to concurrent requests during a slow analysis -- verified live
that \texttt{GET /health} answered in under 3\,ms while a 30-second-plus
area analysis was simultaneously in flight, because the blocking analysis
pipeline is explicitly offloaded to a thread pool rather than run inline
on the async event loop.

\textbf{Availability/Resilience}: the system tolerates the failure of any
single external dependency without failing the whole request -- every
external call degrades to an explicit ``unavailable'' status plus a
warning message rather than raising an exception. This is not a
theoretical concern: every external dependency (Open-Meteo, all three
elevation providers, and Overpass) was observed to fail for real during
development, including DNS resolution failures, connection resets, and
confirmed rate-limiting responses (\S\ref{sec:discussion}).

\textbf{Usability}: single-page web app; dark theme by default with a
persistent light/dark switch; both input modes accessible from one
screen; explicit loading state with a time-expectation message; explicit,
human-readable error banners rather than a blank or broken UI on failure.

\textbf{Security}: no authentication is implemented -- a deliberate scope
decision appropriate for a single-team demo deployment, not a production
multi-tenant system.

\textbf{Browser support}: tested against Chromium-based browsers via
automated browser verification during development; no browser-specific
code paths.

%% =====================================================================
\section{System Architecture and High-Level Design}
\label{sec:architecture}

The system has three tiers plus five external data sources. The client
tier is a React single-page application that either uploads a KML/KMZ
file or submits a drawn polygon as JSON over a CORS-enabled REST API, with
no authentication layer between them. The server tier is a layered FastAPI
application: \texttt{api/} (HTTP route handlers) calls into
\texttt{services/} (pipeline orchestration), which calls
\texttt{algorithms/} (pure computational routines) and \texttt{providers/}
(external API integrations), operating on \texttt{models/} (dataclasses
and Pydantic schemas). Each layer is independently testable. There is no
database tier (\S\ref{sec:implementation}); results are held in an
in-memory store for the life of the server process. Five independent
external providers sit behind the services layer -- Open-Meteo for
rainfall, a three-tier elevation cascade (OpenZenith, then Open-Elevation,
then OpenTopoData), and the Overpass API for OpenStreetMap land-use data
-- each wrapped in its own resilience layer (bounded timeout, bounded
retries, and graceful degradation to an explicit ``unavailable'' result
rather than an exception).

\subsection{Technology Stack}

\textbf{Backend}: Python, FastAPI, Pydantic / pydantic-settings, Uvicorn.
Geospatial: GeoPandas, Shapely, pyproj, Fiona, rasterio. Numerical: NumPy,
SciPy. KML/XML: lxml. HTTP client: httpx. Testing: pytest and
pytest-asyncio (70 tests).

\textbf{Frontend}: React 19, Vite, Leaflet with \texttt{react-leaflet} and
\texttt{leaflet-draw} (pinned to a known-good version to avoid an upstream
bug in later releases), plain CSS with CSS-variable theming. No
TypeScript, no state-management library, no UI framework -- a deliberate
minimalism choice matching the assignment's own emphasis on a clean,
functional interface over a heavy frontend stack.

\textbf{External APIs, all free and requiring no API key}: Open-Meteo
Historical Weather API for rainfall; OpenZenith as the primary elevation
provider for the map-drawn-area path (the assignment brief's own suggested
elevation API); Open-Elevation and OpenTopoData as a two-tier elevation
fallback, tried only for points OpenZenith could not provide; the
Overpass API over OpenStreetMap data for land-use exclusion geometries.

\textbf{Deviation from the suggested stack}: the brief suggested Flask or
FastAPI; FastAPI was chosen for its async support, automatic OpenAPI
documentation, and Pydantic validation, which fit the geospatial and
numerical workload well. No relational database (PostgreSQL with PostGIS
was considered during planning) was used in the final system -- the
in-memory store's own module docstring documents the intended upgrade
path as an isolated, additive change rather than a rearchitecture, which
was judged sufficient for this assignment's scope.

%% =====================================================================
\section{Methodology}
\label{sec:methodology}

\subsection{Terrain and Elevation Analysis}
Two independent ways obtain a Digital Elevation Model, feeding the same
downstream pipeline. From an uploaded KML/KMZ contour map, contour lines
are parsed, sampled along each line, projected into an auto-detected UTM
zone, and interpolated onto a regular grid via \texttt{scipy}'s
\texttt{griddata} (linear interpolation with a nearest-neighbour pass to
fill edge gaps). From a user-drawn map area, the polygon's bounding box
(padded 10\% on each side to avoid flow-routing edge artifacts) defines
the same grid specification, but elevation for every grid cell is fetched
from a live elevation API instead of interpolated: OpenZenith is tried
first (up to 2000 points per batch request), falling back per-point to
Open-Elevation and then OpenTopoData for anything OpenZenith could not
provide. A hard 45-second wall-clock budget bounds the whole fetch
regardless of grid size; any points not fetched in time are filled via
nearest-neighbour interpolation from whatever was fetched, with a warning
reported to the caller. Because a candidate could otherwise be generated
inside that 10\% padding margin -- outside the area the user actually
selected -- an additional hard filter rejects any accepted candidate whose
point falls outside the literal drawn polygon, while catchment
\emph{boundaries} are left unrestricted, since a real drainage basin is
not bounded by an arbitrary rectangle.

Both paths then run identical hydrological conditioning: Priority-Flood
depression filling (Barnes et al., 2014), a min-heap-based algorithm that
raises artificial sinks to their spill elevation, guaranteeing continuous
downhill drainage, and slope computation via a finite-difference gradient.

\subsection{Catchment Area Delineation}
For every cell, D8 flow direction chooses the steepest downhill neighbour
among 8 candidates (distance-weighted elevation drop). Flow accumulation
treats the resulting direction grid as a directed acyclic graph and
processes it in topological order (Kahn's algorithm) to count upstream
contributing cells per cell. For a given candidate pour point, an upstream
breadth-first search on the inverted flow graph collects every cell that
drains to it; the resulting mask is vectorized into a GeoJSON polygon and
its area computed in square kilometres.

Candidate generation draws seeds from two independent signals --
depression seeds (cells raised during depression filling by more than
10\,cm) and flow-convergence seeds (top-1-percentile flow accumulation) --
merges and de-clusters them via spatial non-maximum suppression, then
applies hard filters in sequence before scoring: slope too steep (default
threshold 15$^\circ$), outside the user's drawn selection (map-drawn-area
path only), land-use constraint, and catchment too small (default
threshold 0.005\,km\textsuperscript{2}). The land-use constraint checks
the candidate's point location against a unioned, buffered set of
OpenStreetMap building, road, waterway, water-body, and power-line
geometries fetched live via the Overpass API; a candidate inside any
buffer is rejected with a specific, human-readable reason. Surviving
candidates are scored on five min-max-normalized factors (catchment area,
flow accumulation, slope, relief, and depression depth, weighted 35\%,
25\%, 20\%, 10\%, and 10\% respectively) and ranked.

A stated limitation: the land-use filter checks the candidate point, not
its full catchment polygon -- a catchment naturally spans roads and
rivers upstream, which is hydrologically correct and intentionally not
filtered, but can look like an oversight to a first-time viewer of the
map. Only simple OpenStreetMap ``way'' geometries are handled; complex
multipolygon relations are skipped as a documented simplification. DEM
resolution (30\,m default, configurable) bounds the spatial precision of
every downstream computation.

\subsection{Rainfall Data Integration}
Open-Meteo's Historical Weather API is queried for daily precipitation at
the recommended candidate's coordinates (or the area's centroid if none is
accepted), for a configurable historical window (2015 through the last
complete calendar year by default). Daily values are aggregated into an
annual average, a per-month average across all covered years, and a
monsoon (June--September) seasonal average, feeding the Rational Method
runoff estimate: annual runoff (m\textsuperscript{3}) equals the runoff
coefficient times annual rainfall (m) times catchment area
(m\textsuperscript{2}), with the coefficient defaulting to 0.35
(documented as representing mixed agricultural/semi-pervious land, pending
a future soil- or land-cover-based refinement). Planned pond storage is
80\% of annual runoff by default, with a default depth of 3.0\,m (2.0\,m
on candidates with slope above 5$^\circ$) and surface area derived as
storage divided by depth.

%% =====================================================================
\section{Implementation}
\label{sec:implementation}

\subsection{Backend and API Design}

\begin{table}[h]
  \caption{REST API endpoints}
  \label{tab:endpoints}
  \begin{tabular}{p{5.2cm}p{2.0cm}}
    \toprule
    Method \& Path & Purpose \\
    \midrule
    \texttt{POST /api/v1/analyzeContour} & Upload a KML/KMZ file $\to$ full analysis \\
    \texttt{POST /api/v1/analyzeArea} & Submit a drawn polygon $\to$ full analysis \\
    \texttt{GET /api/v1/analysis/\{id\}} & Retrieve a prior analysis \\
    \texttt{GET /api/v1/health} & Service health check \\
    \bottomrule
  \end{tabular}
\end{table}

No authentication is implemented (\S\ref{sec:requirements}). Every
external call (Open-Meteo, all three elevation providers, Overpass) is
wrapped to catch timeout, HTTP-status, connection, and malformed-payload
errors and convert them into a provider-specific ``unavailable'' result
with a human-readable message, rather than propagating an exception -- the
request still returns \texttt{200 OK} with the rest of the analysis intact
plus a warnings list. A same-second retry is only attempted for failures
that can plausibly succeed on retry (timeouts, server errors); a
DNS/connection failure is not retried, since a same-second retry of a
broken connection essentially never succeeds and was found, in a real
production log, to needlessly double the time lost per failure. Internal
and validation errors use a consistent JSON error shape across all
endpoints, enforced by a custom exception handler.

\subsection{Frontend and Visualization}

A React and Leaflet single-page application renders map layers as GeoJSON
overlays: contour lines (upload mode only), candidate points colour-coded
by status with a distinct marker for the top-ranked candidate, catchment
boundary polygons, and a recommended-location marker, each with a popup
showing its own statistics. The base map is standard OpenStreetMap raster
tiles, with a CSS colour-invert filter applied for the dark theme rather
than switching tile source (chosen after verifying that alternative
``dark'' basemap tile sets have real rural-coverage gaps for the target
region), plus a satellite imagery option. The results panel surfaces
status, warnings, the recommended pond's location, catchment area,
expected annual water volume, and planned storage -- the fields the
assignment brief specifically requires -- alongside supporting
rainfall/runoff detail and a collapsible full candidate table.

\subsection{Database and Storage}
There is no database. Analysis results are held in an in-memory store
keyed by a UUID analysis identifier and are lost on server restart. This
was a deliberate scope decision for a single-session assignment
deployment; the storage module's own docstring documents the upgrade path
to a persistent, spatially-indexed database as an isolated, additive
change that does not require altering the service or API layers.

%% =====================================================================
\section{CSD Themes and Topics Applied in the Project}
\label{sec:csd-themes}

\begin{table*}[t]
  \small
  \renewcommand{\arraystretch}{0.85}
  \caption{Mapping of core CSD themes/topics to their use in this project}
  \label{tab:csd-themes}
  \begin{tabular}{p{2.5cm}p{5.4cm}p{6.4cm}}
    \toprule
    CSD Theme/Topic & Where used in the project & Justification / design rationale \\
    \midrule
    API Design (REST) & Four versioned endpoints under \texttt{/api/v1}, auto-generated OpenAPI docs & Consistent JSON contract, self-documenting for frontend integration; versioning leaves room for a future v2 \\
    Caching & Three in-memory caches: rainfall, elevation, land-use, each keyed by rounded coordinates & Avoids redundant calls to slow/rate-limited public APIs; only \emph{successful} results are cached -- caching a failure was found, via live testing, to permanently disable that feature until process restart \\
    Concurrency / Async & Elevation batches fetched in a thread pool; the blocking pipeline is offloaded from the event loop; land-use is fetched concurrently with elevation & Cuts worst-case fetch time; keeps \texttt{/health} responsive during a slow analysis (a real freeze bug was fixed this way); overlaps two calls that otherwise stacked toward the client timeout \\
    Error Handling / Resilience & Every external call (5 providers) degrades to ``unavailable'' plus a warning instead of raising; elevation and land-use each have a fallback/deadline & All 5 dependencies failed for real during development -- DNS errors, confirmed rate-limiting, connection resets -- the normal operating condition on the deployment network, not a hypothetical \\
    Algorithms \& Complexity & Priority-Flood depression fill, D8 flow direction, flow accumulation (topological sort), watershed BFS, spatial NMS & Chosen over naive/quadratic alternatives since the DEM can be tens of thousands of cells even for a small village \\
    Design Patterns & Interchangeable provider modules behind one function interface; a single orchestrator coordinates independently-testable stages & New providers can be swapped without touching the pipeline; the orchestrator is the only component that knows the full pipeline order \\
    Testing Strategy & 70 automated pytest tests plus live browser-driven verification for the frontend & Confirms resilience behaviour, not just the happy path -- several real bugs were caught by live testing, not mocks \\
    Version Control & Incremental commits per fix/feature, messages documenting root cause and verification & Keeps history a readable record of what was found and why \\
    Load Balancing & Not implemented & Single-instance deployment; out of scope at this scale \\
    DB Indexing / Query Opt. & Not applicable & No persistent database exists (\S\ref{sec:implementation}) \\
    Microservices vs.\ Monolith & Monolith (single FastAPI app) & A distributed split adds operational complexity with no benefit here; the layered structure already separates concerns \\
    Authentication & Not implemented & Explicit scope decision; needed before any public multi-user deployment \\
    Containerization & Not yet done & Open item, see \S\ref{sec:discussion} \\
    \bottomrule
  \end{tabular}
\end{table*}

\emph{Error Handling and Resilience} and \emph{Concurrency} were the
dominant engineering concerns of this project in practice, not just
checklist items. Every external API this system depends on was observed
to genuinely fail during development -- not as a hypothetical edge case,
but repeatedly, in real logs, on the actual deployment network, including
DNS resolution failures, confirmed rate-limiting from at least two
independent providers (one surfaced directly as an HTTP 429 response, the
other traced to a Cloudflare-level ``origin rate limited'' response after
only a handful of rapid requests), and connection resets. The system's
core design principle -- every external call degrades to an explicit
``unavailable'' status plus a warning rather than failing the whole
request -- was validated against these real failures, not just
unit-test mocks. Concurrency work followed directly from a bug this
principle surfaced: a single slow, degraded external-API call was found
to freeze the entire server, because the request handler ran inline on
the async event loop. Fixing it, by offloading to a thread pool, was
necessary specifically because of how failure-prone the external
dependencies turned out to be in practice.

%% =====================================================================
\section{Results and Evaluation}
\label{sec:results}

\begin{table}[h]
  \caption{Representative example, KML upload path (IIT Bhilai contour dataset)}
  \label{tab:results-kml}
  \begin{tabular}{p{4.4cm}p{2.8cm}}
    \toprule
    Metric & Value \\
    \midrule
    Contour lines parsed & 1{,}355 (267.0--298.0\,m) \\
    DEM grid & 106$\times$130 @ 30\,m \\
    Average slope & 1.9$^\circ$ \\
    Candidate seeds (raw $\to$ after NMS) & 418 $\to$ 187 \\
    Candidates scored (top score) & 15 (82.8) \\
    Accepted / rejected (this run) & 5 / 0 \\
    Recommended catchment area & 0.0450\,km\textsuperscript{2} \\
    Historical rainfall (annual avg.) & 1{,}385.4\,mm \\
    Estimated annual runoff & 21{,}820.05\,m\textsuperscript{3} \\
    Planned pond storage & 17{,}456.04\,m\textsuperscript{3} \\
    Recommended depth / surface area & 3.0\,m / 5{,}818.7\,m\textsuperscript{2} \\
    End-to-end processing time & 1.17--7.22\,s \\
    \bottomrule
  \end{tabular}
\end{table}

The identical dataset, analyzed again on a run where Overpass succeeded
(analysis \texttt{62376cca}), rejected 10 of 15 candidates with a
land-use-constraint status against real OpenStreetMap data, leaving the
same top-5 candidates and identical scores as Table~\ref{tab:results-kml}
-- confirming the terrain/scoring pipeline is fully deterministic; only
the external API outcome varies run to run. That same run is also a
clean, real illustration of this project's central resilience claim:
Overpass succeeded while Open-Meteo failed in the same request, and the
response still came back a complete, correctly-filtered candidate list
with \texttt{status: "success"} -- only the rainfall-dependent fields
were null, with an explicit warning explaining why. Each external
dependency degrades independently; a failure in one never blocks a
result the others can still produce.

A second representative run on a $\sim$0.6\,km\textsuperscript{2}
user-drawn polygon over the same region fetched 675 elevation samples
(200 requiring nearest-neighbour fill after provider retries), generated
11 candidates, rejected 6 with a land-use-constraint status against real
OpenStreetMap data (23 roads, 1 waterway, 2 water bodies, and 1 power line
found in the bounding box), and accepted 5. Verified visually in the
browser: rejected candidate markers line up directly along real road
geometry on the satellite basemap, and the recommended marker sits in
open space away from both the river and all roads. No published
ground-truth figures were available to compare against; this is noted
honestly as a limitation rather than fabricated.

\subsection{Performance}

\begin{table}[h]
  \caption{Observed operation times}
  \label{tab:performance}
  \begin{tabular}{p{5.6cm}p{2.2cm}}
    \toprule
    Operation & Observed time \\
    \midrule
    KML-upload analysis, end-to-end & 1.17--7.22\,s \\
    Map-drawn-area, real fetch, no fallback needed & 1--2\,s \\
    Map-drawn-area, cached elevation & 0.02--0.06\,s \\
    Map-drawn-area, providers degraded (worst case) & bounded $\sim$75--90\,s \\
    \texttt{GET /health} during a 30\,s+ analysis & $<$3\,ms \\
    Overpass land-use query, healthy mirror & 1--2\,s \\
    \bottomrule
  \end{tabular}
\end{table}

No formal load testing (many concurrent users) was performed; this is
noted as a limitation, not claimed as tested.

%% =====================================================================
\section{Discussion and Limitations}
\label{sec:discussion}

\textbf{External API reliability was the single biggest practical risk to
this system's usability}, not any algorithmic weakness. During
development, and especially when testing on the actual grading/lab
machine rather than an isolated development sandbox, every external
dependency showed real transient failures -- DNS resolution failures,
connection resets, ``network unreachable'' errors, and gateway timeouts,
sometimes several within a single request. This confirmed the
unreliability is a property of that network's egress in general, not any
one provider being uniquely bad, and directly validated the project's
resilience-first design rather than being a hypothetical concern. Direct
testing also confirmed that two of the five providers enforce real,
mostly undocumented rate limits: one returned an explicit HTTP 429
response under load, and another began returning a Cloudflare-level
``origin rate limited'' response after only a handful of requests in
quick succession from a single network. With an entire class concurrently
exercising the same free services, likely from one shared institutional
IP address, these limits are easy to trip continuously, independent of
any code-level bug.

Live testing surfaced several genuine bugs that no unit test caught on
its own. When an elevation provider failed a large batch late, after
already spending much of the shared time budget on its own retries, the
fallback chain re-chunked that batch into many small sub-batches but
originally fetched them sequentially, so only the first one or two could
complete before the deadline -- a batch that was genuinely recoverable
ended up mostly missing, occasionally triggering a false
``elevation unavailable'' error; this was fixed by fetching sub-batches
concurrently. Separately, the elevation point cache was found to store a
failed lookup exactly as it would store a successful one, so a single bad
network moment permanently disabled elevation for that area until the
server process was restarted, even after the network fully recovered;
this was fixed by only caching successful values. A third issue was
architectural rather than a bug: sequential external calls (elevation,
land-use, rainfall) were found, from a real log, to stack close enough to
the frontend's own request timeout that a slightly worse network moment
could cause the client to abandon a request the backend was about to
complete successfully; this was addressed by bounding the land-use
fetch's own worst case and launching it concurrently with the elevation
fetch, since both depend only on the drawn area's bounding box.

Further limitations, stated honestly rather than omitted: the land-use
filter checks a candidate's point, not its full catchment polygon, so a
recommended pond's catchment can legitimately cross a road or river
upstream -- hydrologically correct, but visually surprising to a
first-time viewer; the runoff coefficient is a single configurable
constant rather than derived from soil type or land cover, a real
accuracy limitation flagged explicitly in the UI's assumptions section;
there is no database, so results do not survive a server restart; no
formal load testing was performed; and DEM resolution (30\,m default)
bounds the spatial precision of every downstream computation. A
third-party dependency risk was also found and worked around: the map's
polygon-drawing library has a known upstream bug in its latest release
that breaks rectangle drawing under modern strict-mode bundlers, worked
around by pinning to an earlier, unaffected release.

%% =====================================================================
\section{AI Tool Usage Declaration}
\label{sec:ai-usage}

Claude (Anthropic's Claude Code CLI) assisted with implementation, debugging,
and report drafting; all changes were reviewed and tested by the author.

%% =====================================================================
\appendix

\section{Source Code and Repository}
\textbf{Repository}: \url{https://github.com/SandstromPL/AI-based-Village-Pond-Planning-System-}.
\textbf{Deployed frontend}: \url{http://10.1.75.53:3279/}.
\textbf{Demo video}: \url{https://youtu.be/MARQ1Q1CHxE}.

Folder structure: \texttt{backend/app/} is split into \texttt{api/}
(HTTP route handlers), \texttt{services/} (pipeline orchestration, one
service per concern), \texttt{algorithms/} (pure computational routines:
depression filling, D8 flow, flow accumulation, watershed, scoring,
slope, interpolation), \texttt{providers/} (external API integrations for
elevation and land-use), \texttt{models/} (dataclasses and Pydantic
schemas), and \texttt{utils/} (geo helpers, GeoJSON, in-memory storage);
\texttt{backend/tests/} holds 70 pytest tests. \texttt{frontend/src/}
holds \texttt{components/} (map view, results panel, upload/draw
controls, loading and error UI), \texttt{utils/} (display formatting and
defensive GeoJSON filtering), and the top-level API client, theme, and
application entry point.

\end{document}
\endinput
